\chapter{Realizzazione di classi in Java}
\label{chap:realizzazione-classi-java}

Nei \cref{chap:introduzione-java,chap:collaudo-tipi-input-java}
sono state utilizzate classi già disponibili. Per realizzarne una
occorre definirne operazioni, stato e comportamento, distinguendo
interfaccia e implementazione come per gli ADT
(\cref{sec:fi-adt-strutture}). Si utilizzerà come modello
una classe per la gestione di un conto bancario.\footnote{Materiale di riferimento: S. Canazza,
\emph{Realizzare classi}, \nolinkurl{2026_FI_09_canazza_Java_RealizzareClassi.pdf}.}

\section{Astrazione e approccio black-box}
\label{sec:fi-java-classi-black-box}

L'\emph{approccio black-box}, o a scatola nera, consente di
utilizzare un componente conoscendone le operazioni e i risultati
attesi, senza esaminarne la realizzazione interna. Analogamente,
un programma utilizza un oggetto attraverso i suoi metodi,
indipendentemente dal codice che li implementa.

\begin{definition}[Astrazione]
\label{def:fi-java-classi-astrazione}
L'\emph{astrazione} è il processo mediante il quale si costruisce
un modello di un'entità, individuandone le caratteristiche e i
comportamenti rilevanti rispetto allo scopo del programma e
trascurando i dettagli non necessari.
\end{definition}

Per rappresentare un conto bancario, per esempio, possono essere
sufficienti la possibilità di consultare il saldo, versare e
prelevare denaro. Non è necessario riprodurre ogni aspetto di una
banca reale.

\begin{definition}[Incapsulamento]
\label{def:fi-java-classi-incapsulamento}
L'\emph{incapsulamento} è il principio della programmazione
orientata agli oggetti secondo cui lo stato interno e le operazioni
che lo gestiscono sono riuniti in una classe, limitando l'accesso
diretto alla rappresentazione interna attraverso un'interfaccia
controllata.
\end{definition}

Astrazione e incapsulamento sono collegati, ma non equivalenti:
la prima guida la scelta del modello, mentre il secondo disciplina
l'accesso alla sua realizzazione. La
\cref{fig:fi-java-classi-black-box} rappresenta questa separazione.

\begin{figure}[htbp]
  \centering
  \begin{tikzpicture}[unipd/diagram,x=1cm,y=1cm]
  \node[unipd/node,minimum width=2.2cm,minimum height=1.15cm,
    align=center] (client) at (-3.7,0)
    {Programma\\utilizzatore};
  \node[unipd/node,minimum width=2.9cm,minimum height=1.15cm,
    align=center] (api) at (0,0)
    {Interfaccia pubblica\\\texttt{deposit(), withdraw()}\\\texttt{getBalance()}};
  \node[unipd/node,minimum width=2.35cm,minimum height=1.15cm,
    align=center] (state) at (3.9,0)
    {Stato interno\\\texttt{private balance}};
  \draw[unipd/directed edge] (client.east) -- (api.west);
  \draw[unipd/directed edge] (api.east) -- (state.west);
\end{tikzpicture}

  \caption{Utilizzo di un oggetto attraverso l'interfaccia pubblica,
  senza accesso diretto alla rappresentazione interna.}
  \label{fig:fi-java-classi-black-box}
\end{figure}

\section{Interfaccia pubblica e definizione dei metodi}
\label{sec:fi-java-classi-metodi}

Si consideri una classe \texttt{BankAccount} destinata a
rappresentare un conto corrente. Le operazioni essenziali del modello
sono il versamento di una somma, il prelievo e la consultazione del
saldo. Le operazioni possono essere descritte schematicamente mediante
le seguenti dichiarazioni, che ne omettono il corpo:

\begin{unipdcode}[language=Java]
public void deposit(double amount);
public void withdraw(double amount);
public double getBalance();
\end{unipdcode}

\clearpage

Le dichiarazioni riportate descrivono le operazioni disponibili,
ma non costituiscono una definizione completa dei metodi in una
classe concreta. I metodi
\texttt{deposit()} e \texttt{withdraw()} ricevono l'importo da
elaborare e non restituiscono un valore; \texttt{getBalance()}
restituisce invece un \texttt{double} e non richiede argomenti.

\begin{definition}[Interfaccia pubblica]
\label{def:fi-java-classi-interfaccia}
L'\emph{interfaccia pubblica} di una classe è l'insieme dei
membri accessibili dall'esterno che definiscono le modalità
di interazione con la classe e i suoi oggetti, specificandone
le operazioni disponibili e il comportamento osservabile,
indipendentemente dai dettagli della loro implementazione.
\end{definition}

In particolare, i costruttori e i metodi pubblici stabiliscono
come creare e utilizzare gli oggetti, quali argomenti fornire,
quali risultati attendersi e quali condizioni rispettare.
L'interfaccia pubblica non deve esporre tutte le operazioni
tecnicamente possibili sui dati, ma soltanto quelle previste
dal modello. Per esempio, conoscere il saldo di un conto non
implica poterlo impostare arbitrariamente.

La specifica dei metodi deve consentirne l'uso e il collaudo
senza richiederne il codice.

\subsection{Intestazione, corpo e firma di un metodo}
\label{sec:fi-java-classi-firma}

La definizione di un metodo è composta da un'\emph{intestazione}
e da un \emph{corpo}, racchiuso fra parentesi graffe. Per esempio,

\begin{unipdcode}[language=Java]
public void deposit(double amount) {
    balance = balance + amount;
}
\end{unipdcode}

nell'intestazione \texttt{public} indica l'accessibilità,
\texttt{void} il tipo di ritorno, \texttt{deposit} il nome del
metodo e \texttt{double amount} il parametro dichiarato. Il corpo
contiene le istruzioni eseguite quando il metodo viene invocato.
Se non è necessario ricevere argomenti, l'elenco dei parametri
rimane vuoto, come in \texttt{getBalance()}.

La \emph{firma} di un metodo Java, in senso tecnico, comprende il
nome e i tipi dei parametri, considerandone anche l'ordine; non
comprende il tipo restituito né i modificatori di accesso. Questa
distinzione è importante nella risoluzione dei metodi sovraccarichi,
trattata nella \cref{sec:fi-java-classi-overloading}. Per gli
esempi di invocazione si rimanda anche alla
\cref{sec:fi-java-invocazione-metodi}.

\begin{definition}[Metodi di accesso e modificatori]
\label{def:fi-java-classi-accessori}
Un \emph{metodo di accesso} (\emph{accessor}) osserva lo stato di
un oggetto senza modificarlo. Un \emph{metodo modificatore}
(\emph{mutator}) può invece alterarne lo stato.
\end{definition}

Nella classe proposta, \texttt{getBalance()} è un metodo di
accesso, mentre \texttt{deposit()} e \texttt{withdraw()} sono
modificatori. Anche \texttt{length()} della classe \texttt{String}
è un metodo di accesso; \texttt{translate()} di
\texttt{Rectangle} modifica invece l'oggetto ricevente
(\cref{sec:fi-java-riferimenti}). La distinzione riguarda gli
effetti sullo stato, non il tipo restituito: un metodo modificatore
può anche restituire un valore.

\clearpage

\subsection{Parametri espliciti e valore restituito}
\label{sec:fi-java-classi-parametri}

La chiamata a un metodo di istanza identifica un oggetto
\emph{ricevente} e fornisce gli eventuali argomenti. Per esempio,

\begin{unipdcode}[language=Java]
account.deposit(500);
double balance = account.getBalance();
\end{unipdcode}

nella prima istruzione il riferimento \texttt{account} individua
l'oggetto sul quale operare e \texttt{500} è l'argomento associato
al parametro esplicito \texttt{amount}. Il ricevente costituisce il
\emph{parametro implicito} del metodo: l'istruzione che modifica
\texttt{balance} agisce sul campo dell'oggetto indicato da
\texttt{account}, non su quello di tutti gli oggetti della classe.
All'interno di un metodo di istanza, il ricevente può essere
indicato esplicitamente con \texttt{this}, come in
\texttt{this.balance}.

L'istruzione \texttt{return} termina l'esecuzione del metodo e
restituisce il controllo al chiamante. In un metodo che produce
un risultato, \texttt{return espressione;} ne specifica il valore,
compatibile con il tipo dichiarato; un metodo \texttt{void} può
usare \texttt{return;} per terminare anticipatamente oppure
raggiungere la fine del corpo. Un metodo con un tipo di ritorno
diverso da \texttt{void} deve invece restituire un valore lungo
ogni percorso che termina normalmente.

\section{Costruttori e inizializzazione degli oggetti}
\label{sec:fi-java-classi-costruttori}

Per usare un oggetto occorre prima crearlo e inizializzarlo. In Java,
l'operatore \texttt{new} alloca un nuovo oggetto, ne determina
l'inizializzazione e produce un riferimento al risultato. La
costruzione \texttt{new BankAccount()} può essere usata in un
assegnamento della forma

\begin{unipdcode}[language=Java]
BankAccount account = new BankAccount();
\end{unipdcode}

Il riferimento viene assegnato alla variabile \texttt{account};
la variabile e l'oggetto non sono la stessa entità, come chiarito
nella \cref{sec:fi-java-riferimenti}.

\begin{definition}[Costruttore]
\label{def:fi-java-classi-costruttore}
Un \emph{costruttore} è una dichiarazione speciale di una classe,
eseguita durante la creazione di una sua istanza per inizializzarne
lo stato. Ha lo stesso nome della classe e non dichiara alcun tipo
di ritorno, neppure \texttt{void}.
\end{definition}

Per \texttt{BankAccount} si possono prevedere due modi di
inizializzazione: un saldo iniziale nullo oppure un saldo indicato
dal chiamante. I relativi costruttori, da collocare nel corpo
della classe, sono

\begin{unipdcode}[language=Java]
public BankAccount() {
    balance = 0;
}

public BankAccount(double initialBalance) {
    balance = initialBalance;
}
\end{unipdcode}

Le istruzioni \texttt{new BankAccount()} e
\texttt{new BankAccount(500)} eseguono costruttori differenti.
L'operatore \texttt{new} produce il riferimento all'oggetto:
non è il costruttore a restituirlo. La dichiarazione
\texttt{public void BankAccount()} non definirebbe invece un
costruttore, ma un metodo ordinario con lo stesso nome della
classe; una sua invocazione non inizializzerebbe automaticamente
un nuovo oggetto.

\subsection{Costruttore predefinito e valori iniziali}
\label{sec:fi-java-classi-costruttore-predefinito}

Se una classe non dichiara \emph{alcun} costruttore, il
compilatore ne fornisce uno \emph{predefinito} senza parametri,
con accessibilità coerente con quella della classe. Quando invece
viene dichiarato almeno un costruttore, quello predefinito non
viene generato: definendo soltanto
\texttt{BankAccount(double initialBalance)}, la chiamata
\texttt{new BankAccount()} non sarebbe valida.

I campi di istanza e quelli statici ricevono valori iniziali
predefiniti anche quando non sono assegnati esplicitamente,
secondo le regole della \cref{tab:fi-java-valori-predefiniti}.

\begin{table}[htbp]
  \centering
  \begin{tabular}{ll}
    \toprule
    \textbf{Tipo} & \textbf{Valore iniziale} \\
    \midrule
    \texttt{byte}, \texttt{short}, \texttt{int}, \texttt{long}
      & \texttt{0} \\
    \texttt{float} & \texttt{0.0f} \\
    \texttt{double} & \texttt{0.0d} \\
    \texttt{boolean} & \texttt{false} \\
    \texttt{char} & \texttt{\textquotesingle\textbackslash u0000\textquotesingle} \\
    Tipi riferimento & \texttt{null} \\
    \bottomrule
  \end{tabular}
  \caption{Valori iniziali predefiniti dei campi in Java.}
  \label{tab:fi-java-valori-predefiniti}
\end{table}

Questa regola non va confusa con l'inizializzazione delle variabili
locali: il compilatore ne impedisce la lettura se non può stabilire
che abbiano ricevuto un valore lungo tutti i percorsi di esecuzione
possibili.

\section{Sovraccarico dei metodi e dei costruttori}
\label{sec:fi-java-classi-overloading}

I due costruttori di \texttt{BankAccount} hanno lo stesso nome,
ma elenchi di parametri differenti. Questa possibilità prende
il nome di \emph{sovraccarico} (\emph{overloading}) ed è prevista
anche per i metodi ordinari.

\begin{definition}[Overloading]
\label{def:fi-java-classi-overloading}
Si ha \emph{sovraccarico} quando una classe dichiara più metodi
con lo stesso nome, oppure più costruttori, distinguibili per
numero, tipo o ordine dei parametri. Il compilatore risolve
l'invocazione individuando la dichiarazione applicabile agli
argomenti forniti secondo le regole del linguaggio.
\end{definition}

Il tipo del risultato non distingue due firme: non è quindi
possibile definire due metodi aventi lo stesso nome e lo stesso
elenco di parametri, cambiando soltanto il tipo restituito.
Se nessuna dichiarazione è applicabile, la compilazione fallisce,
come accadrebbe con \texttt{new BankAccount("testo")} quando
sono presenti soltanto i due costruttori precedenti. Anche la
presenza di più dichiarazioni applicabili senza una scelta
univoca può provocare un errore di ambiguità.

\begin{example}[Metodi sovraccarichi nella libreria standard]
La classe \texttt{PrintStream}, a cui appartiene
\texttt{System.out}, fornisce versioni di \texttt{println()}
per argomenti di tipi differenti. Per esempio:

\begin{unipdcode}[language=Java]
System.out.println(10);       // Argomento int
System.out.println(10.5);     // Argomento double
System.out.println("Dieci");  // Argomento String
\end{unipdcode}
Il compilatore seleziona il metodo appropriato senza
richiedere nomi differenti per ciascun tipo.
\end{example}

\clearpage

\section{Classi, oggetti e variabili di istanza}
\label{sec:fi-java-classi-stato}

\begin{definition}[Classe e oggetto]
\label{def:fi-java-classi-classe-oggetto}
Una \emph{classe} è un modello che definisce la struttura e
il comportamento comuni a un insieme di oggetti. Un
\emph{oggetto} è un'istanza concreta di una classe, creata
durante l'esecuzione del programma e dotata di un proprio stato.
\end{definition}

\begin{definition}[Stato e variabili di istanza]
\label{def:fi-java-classi-stato}
Lo \emph{stato} di un oggetto è determinato dai valori dei suoi
\emph{campi di istanza}, detti anche \emph{variabili di esemplare}.
Ogni oggetto possiede una propria copia di tali campi, mentre
le definizioni dei metodi di istanza sono comuni alla classe.
\end{definition}

Questa distinzione approfondisce quanto introdotto nella
\cref{sec:fi-java-classi-oggetti}.

Nel conto bancario è sufficiente rappresentare il saldo con
un campo di tipo \texttt{double}, dichiarato all'interno
della classe:

\begin{unipdcode}[language=Java]
private double balance;
\end{unipdcode}

La dichiarazione del campo appartiene alla classe, ma ogni
oggetto \texttt{BankAccount} mantiene un proprio valore di
\texttt{balance}. Si considerino le istruzioni

\begin{unipdcode}[language=Java]
BankAccount account1 = new BankAccount(1000);
BankAccount account2 = new BankAccount(2000);
\end{unipdcode}

Le due istruzioni creano istanze distinte, come illustrato nella
\cref{fig:fi-java-classi-istanze}. Una modifica del primo
saldo non cambia quello del secondo.

\begin{figure}[htbp]
  \centering
  \begin{tikzpicture}[unipd/diagram,x=1cm,y=1cm]
  \node[unipd/node,minimum width=1.85cm] (a1) at (-3.4,0.85)
    {\texttt{account1}};
  \node[unipd/node,minimum width=1.85cm] (a2) at (-3.4,-0.85)
    {\texttt{account2}};
  \node[unipd/node,minimum width=3.55cm,minimum height=0.85cm,
    align=center] (o1) at (1.25,0.85)
    {\texttt{BankAccount}\\\texttt{balance = 1000.0}};
  \node[unipd/node,minimum width=3.55cm,minimum height=0.85cm,
    align=center] (o2) at (1.25,-0.85)
    {\texttt{BankAccount}\\\texttt{balance = 2000.0}};
  \draw[unipd/directed edge] (a1.east) -- (o1.west);
  \draw[unipd/directed edge] (a2.east) -- (o2.west);
\end{tikzpicture}

  \caption{Due variabili riferite a istanze distinte, ciascuna
  con il proprio campo \texttt{balance}.}
  \label{fig:fi-java-classi-istanze}
\end{figure}

L'assegnamento \texttt{BankAccount alias = account1;}
non crea invece una terza istanza: \texttt{alias} e
\texttt{account1} si riferiscono allo stesso oggetto.
La differenza fra copia di un riferimento e creazione
di un oggetto distinto è già mostrata nella
\cref{fig:fi-java-alias-copia}.

\subsection{Riferimenti nulli}
\label{sec:fi-java-classi-null}

Una variabile di tipo riferimento può contenere
\texttt{null}, valore che indica l'assenza di un
oggetto. Dopo l'istruzione

\begin{unipdcode}[language=Java]
BankAccount account = null;
\end{unipdcode}

la variabile è inizializzata, ma non si riferisce ad
alcun conto bancario. Invocare su di essa un metodo
di istanza, come \texttt{account.getBalance()}, provoca
una \texttt{NullPointerException}. Se tale eccezione
non viene gestita, l'esecuzione del programma può
interrompersi (\cref{sec:fi-java-errori}).

\begin{example}[Stringa vuota e riferimento nullo]
Si considerino i seguenti riferimenti a oggetti
\texttt{String}:

\begin{unipdcode}[language=Java]
String greeting = "Hello";
String emptyString = "";
String nullString = null;

int a = greeting.length();    // 5
int b = emptyString.length(); // 0
int c = nullString.length();  // NullPointerException
\end{unipdcode}

La stringa vuota è un oggetto valido di lunghezza zero;
\texttt{null} non identifica invece alcun oggetto. La
terza invocazione genera quindi un errore durante
l'esecuzione, non un errore di sintassi.
\end{example}

\section{Incapsulamento e realizzazione di \texttt{BankAccount}}
\label{sec:fi-java-classi-incapsulamento}

Il modificatore \texttt{private} impedisce a codice
esterno alla classe di accedere direttamente al campo
\texttt{balance}. L'espressione \texttt{account.balance},
per esempio, non è consentita in una classe di collaudo,
mentre la chiamata \texttt{account.getBalance()} utilizza
l'interfaccia pubblica prevista dal progettista. I metodi
della stessa classe possono accedere ai campi privati
delle sue istanze.

L'incapsulamento consente di controllare le modifiche
allo stato dell'oggetto. Se un conto non può assumere saldo
negativo, questa proprietà deve essere stabilita come
\emph{invariante}, ossia una condizione che deve valere in ogni
stato ammesso dell'oggetto. I costruttori devono soddisfarla
inizialmente e i metodi modificatori devono preservarla.
Esporre un metodo generico
\texttt{setBalance()} permetterebbe invece di modificare
il saldo senza passare attraverso le operazioni lecite
di versamento e prelievo.

Una prima implementazione di \texttt{BankAccount} è la seguente.
I costruttori inizializzano il saldo, \texttt{deposit()} lo
incrementa, \texttt{withdraw()} lo decrementa e
\texttt{getBalance()} ne restituisce il valore corrente.

\begin{unipdcode}[language=Java]
public class BankAccount {
    private double balance;

    public BankAccount() {
        balance = 0;
    }

    public BankAccount(double initialBalance) {
        balance = initialBalance;
    }

    public void deposit(double amount) {
        balance = balance + amount;
    }

    public void withdraw(double amount) {
        balance = balance - amount;
    }

    public double getBalance() {
        return balance;
    }
}
\end{unipdcode}


\begin{warning}[Limiti dell'implementazione essenziale]
\label{warn:fi-java-classi-bank-limiti}
In questa versione non si controllano gli importi: un
prelievo superiore al saldo può renderlo negativo, mentre
un importo negativo può invertire l'effetto di un'operazione.
Il modificatore \texttt{private} impedisce l'accesso diretto
al campo dall'esterno, ma non garantisce da solo la validità
dello stato. Inoltre, l'uso di \texttt{double} per gli importi
monetari è una semplificazione didattica, soggetta agli
errori di rappresentazione descritti nella
\cref{sec:fi-java-precisione}.
\end{warning}

\subsection{Collaudo della classe}
\label{sec:fi-java-classi-collaudo}

Una classe di dominio come \texttt{BankAccount} non deve
necessariamente contenere un metodo \texttt{main}. Per
collaudarla si può creare una classe separata che costruisca
oggetti, invochi i loro metodi e confronti i risultati
con quelli previsti, secondo il procedimento della
\cref{sec:fi-java-collaudo}.

Il programma seguente parte da un saldo di \(10\,000\) euro
e accredita un interesse annuo del \(5\%\) per due anni,
senza prelievi o ulteriori versamenti. Gli interessi
vengono aggiunti al saldo alla fine di ciascun anno.

\begin{unipdcode}[language=Java]
public class BankAccountTester {
    public static void main(String[] args) {
        BankAccount account = new BankAccount(10000);
        final double RATE = 0.05;

        double interest = account.getBalance() * RATE;
        account.deposit(interest);
        System.out.println("Saldo dopo un anno: " + account.getBalance());

        interest = account.getBalance() * RATE;
        account.deposit(interest);
        System.out.println("Saldo dopo due anni: " + account.getBalance());
    }
}
\end{unipdcode}


Per riprodurre l'esempio, i file \texttt{BankAccount.java}
e \texttt{BankAccountTester.java} devono contenere le
rispettive classi pubbliche e possono essere compilati
e avviati, nella stessa directory e nello stesso pacchetto,
con i comandi

\begin{unipdterminal}
$ javac BankAccount.java BankAccountTester.java
$ java BankAccountTester
Saldo dopo un anno: 10500.0
Saldo dopo due anni: 11025.0
\end{unipdterminal}

Il risultato finale corrisponde alla capitalizzazione composta:
\[
  10\,000(1+0{,}05)^2=11\,025.
\]
Un collaudo completo rispetto al comportamento dichiarato
deve verificare anche la costruzione senza argomenti, i
versamenti, i prelievi e le condizioni limite. I test
possono individuare errori, ma non dimostrano da soli la
correttezza generale dell'implementazione.

\clearpage

\section{Documentazione dell'interfaccia con JavaDoc}
\label{sec:fi-java-classi-javadoc}

La documentazione di una classe dovrebbe descriverne lo
scopo e chiarire il comportamento dei metodi pubblici:
parametri, valori restituiti, condizioni richieste ed
eventuali eccezioni. A questo scopo Java mette a
disposizione \emph{JavaDoc}, uno strumento che genera
pagine HTML a partire da appositi commenti nel sorgente.
Un commento di documentazione inizia con
\lstinline|/**| e termina con \lstinline|*/|, a
differenza dei commenti ordinari introdotti nella
\cref{sec:fi-java-sintassi}.

\begin{example}[Documentazione di un metodo]
Il metodo di accesso al saldo può essere documentato con
un commento JavaDoc:

\begin{unipdcode}[language=Java]
/**
 * Restituisce il saldo corrente del conto.
 *
 * @return saldo attuale
 */
public double getBalance() {
    return balance;
}
\end{unipdcode}
Il commento descrive il risultato osservabile senza
vincolare il codice a una particolare realizzazione
interna del metodo.
\end{example}

I principali \emph{tag} JavaDoc, introdotti dal carattere
\texttt{@}, sono riepilogati nella
\cref{tab:fi-java-classi-javadoc-tag}.

\begin{table}[htbp]
  \centering
  \begin{tabular}{ll}
    \toprule
    \textbf{Tag} & \textbf{Funzione} \\
    \midrule
    \texttt{@param} & descrive un parametro \\
    \texttt{@return} & descrive il valore restituito \\
    \texttt{@throws} & documenta un'eccezione possibile \\
    \texttt{@see} & rimanda a un elemento correlato \\
    \texttt{@author}, \texttt{@version} & indicano autore e versione \\
    \texttt{@since} & indica la versione di introduzione \\
    \texttt{@deprecated} & segnala un elemento obsoleto \\
    \lstinline|{@code ...}| & formatta codice nel testo \\
    \lstinline|{@link ...}| & crea un collegamento a un altro elemento \\
    \bottomrule
  \end{tabular}
  \caption{Principali tag per documentare classi e metodi con JavaDoc.}
  \label{tab:fi-java-classi-javadoc-tag}
\end{table}

\clearpage

\begin{example}[Parametri, risultato ed eccezione]
Il seguente metodo usa \texttt{Math.addExact()} per
sommare due interi. In caso di overflow segnala un errore
(\cref{sec:fi-java-overflow}):
\begin{unipdcode}[language=Java]
/**
 * Calcola la somma esatta di due interi.
 *
 * @param a primo addendo
 * @param b secondo addendo
 * @return somma dei due addendi
 * @throws ArithmeticException se si verifica overflow
 * @see Math#addExact(int,int)
 */
public static int somma(int a, int b) {
    return Math.addExact(a, b);
}
\end{unipdcode}
Il metodo è statico (\cref{sec:fi-java-math}); il commento
ne documenta il contratto, compresa l'eccezione.
\end{example}

I tag \lstinline|{@inheritDoc}| e \texttt{@hidden} permettono,
quando applicabili, di riutilizzare documentazione ereditata o
escludere elementi dalla documentazione generata. Le annotazioni
\texttt{@Deprecated} e \texttt{@Override} appartengono invece
al linguaggio Java, non ai commenti JavaDoc.
Il comando \texttt{javadoc -d docs BankAccount.java} genera
la documentazione HTML nella directory \texttt{docs}.

Nel progetto di una classe conviene quindi procedere
dall'identificazione delle operazioni alla definizione e
documentazione dell'interfaccia, poi individuare i campi,
realizzare costruttori e metodi e infine collaudarli.
La documentazione serve a rendere utilizzabile la classe
anche quando la sua implementazione non è disponibile.

\section{Istruzioni condizionali e confronti}
\label{sec:fi-java-classi-condizioni}

La versione essenziale di \texttt{BankAccount} modifica
il saldo senza eseguire controlli. Per subordinare
un'operazione a una condizione si utilizza l'istruzione
\texttt{if}. A differenza di Python, dove i blocchi
sono individuati dall'indentazione, Java impiega
parentesi tonde per la condizione e, quando necessari,
blocchi delimitati da graffe. Il collegamento con le
istruzioni condizionali già incontrate in Python è
richiamato nella \cref{sec:fi-python-condizioni}.

\subsection{La struttura \texttt{if}/\texttt{else}}
\label{sec:fi-java-classi-if-else}

La condizione di un \texttt{if} deve essere un'espressione
di tipo \texttt{boolean}. Se è vera, viene eseguito il
primo ramo; quando è presente \texttt{else}, il secondo
ramo viene eseguito soltanto se la condizione è falsa:

\begin{unipdcode}[language=Java]
if (condizione) {
    // Istruzioni per il caso vero.
} else {
    // Istruzioni per il caso falso.
}
\end{unipdcode}

La parte \texttt{else} è facoltativa. Se un ramo
contiene un solo enunciato, le graffe possono essere
omesse; è comunque opportuno mantenerle per rendere
esplicito il blocco ed evitare errori quando se ne
aggiungono altre istruzioni. Le graffe delimitano un blocco,
che Java considera un unico enunciato composto anche quando
contiene più istruzioni.

\begin{example}[Prelievo condizionato dal saldo]
Supponendo che l'importo sia finito e non negativo, il metodo
\texttt{withdraw()} può rifiutare un prelievo superiore
al saldo:
\begin{unipdcode}[language=Java]
public void withdraw(double amount) {
    if (amount <= balance) {
        balance = balance - amount;
    } else {
        System.out.println("Conto scoperto");
    }
}
\end{unipdcode}
La modifica viene eseguita soltanto quando la somma
richiesta non supera quella disponibile. Il metodo
non controlla ancora esplicitamente un eventuale
argomento negativo: la sua validità è un requisito
ulteriore da stabilire nella specifica.
\end{example}

Non è equivalente utilizzare due condizioni indipendenti
quando la prima può cambiare le variabili lette dalla
seconda. Si consideri

\begin{unipdcode}[language=Java]
if (amount <= balance) {
    balance = balance - amount;
}
if (amount > balance) {
    System.out.println("Conto scoperto");
}
\end{unipdcode}

Con \texttt{balance = 100} e \texttt{amount = 60}, il
primo controllo riesce e riduce il saldo a \(40\). Il
secondo confronta però \(60\) con il saldo aggiornato,
stampando erroneamente il messaggio di rifiuto. Un unico
\texttt{if}/\texttt{else} realizza invece due alternative
mutuamente esclusive, valutate a partire da una sola
condizione.

\subsection{Operatori relazionali e tipi riferimento}
\label{sec:fi-java-classi-confronti}

Gli operatori relazionali \texttt{<}, \texttt{<=},
\texttt{>} e \texttt{>=} confrontano l'ordine dei valori;
gli operatori di uguaglianza \texttt{==} e \texttt{!=}
verificano invece l'uguaglianza o la disuguaglianza.
L'operatore \texttt{==} è distinto dall'assegnamento
\texttt{=}: il primo produce un \texttt{boolean},
mentre il secondo assegna un valore. Un'espressione numerica non può essere
utilizzata direttamente come condizione, poiché in Java
la condizione di \texttt{if} deve essere booleana.

Per i tipi primitivi numerici, \texttt{==} confronta i
valori dopo le eventuali conversioni numeriche previste.
Per i tipi riferimento, invece, \texttt{==} verifica se
i due riferimenti identificano lo stesso oggetto (oppure
sono entrambi \texttt{null}); non confronta generalmente
i contenuti. Per confrontare il contenuto delle stringhe
si usa invece
\texttt{equals()}, che richiede un riferimento ricevente non nullo.
La differenza fra copia di un riferimento e creazione di un nuovo
oggetto è illustrata nella \cref{sec:fi-java-riferimenti};
il caso del riferimento nullo è discusso nella
\cref{sec:fi-java-classi-null}.

\clearpage

\subsection{Confronti tra numeri in virgola mobile}
\label{sec:fi-java-classi-floating}

La rappresentazione in virgola mobile ha precisione
finita: valori decimali e risultati di calcoli possono
subire approssimazioni
(\cref{sec:fi-virgola-mobile,sec:fi-java-precisione}).
Per esempio, l'espressione \texttt{0.1 + 0.2} non
produce in \texttt{double} lo stesso valore
rappresentato dal letterale \texttt{0.3}. Un confronto
con \texttt{==} restituisce quindi \texttt{false},
benché i due risultati siano matematicamente uguali.

Quando un'applicazione vuole verificare la
\emph{vicinanza numerica} anziché l'uguaglianza esatta,
si può introdurre una tolleranza. Dati due valori
finiti \(x\) e \(y\), il criterio assoluto
\[
  |x-y|\leq\varepsilon_{\mathrm{abs}}
\]
considera accettabile una differenza fissa, mentre il
criterio relativo
\[
  |x-y|\leq\varepsilon_{\mathrm{rel}}
  \max\{|x|,|y|\}
\]
scala la tolleranza con la grandezza dei valori.
Il solo criterio relativo non è però adatto ai
confronti prossimi a zero: se uno dei due valori è
nullo e l'altro è non nullo, non può risultare vero
per \(0<\varepsilon_{\mathrm{rel}}<1\).

Una scelta possibile, per un problema che lo consenta,
è combinare le due tolleranze:
\[
  |x-y|\leq\varepsilon_{\mathrm{abs}}+
  \varepsilon_{\mathrm{rel}}\max\{|x|,|y|\}.
\]
Le due soglie non sono costanti universali: dipendono
dalla precisione richiesta, dalla scala dei dati e
dal significato del confronto. La classe seguente
illustra l'implementazione di tale criterio, escludendo
i valori non finiti.

\begin{unipdcode}[language=Java]
public class Numeric {
    public static boolean approxEqual(double x, double y) {
        final double ABS_TOL = 1E-14;
        final double REL_TOL = 1E-14;
        if (!Double.isFinite(x) || !Double.isFinite(y)) {
            return false;
        }
        double scale = Math.max(Math.abs(x), Math.abs(y));
        return Math.abs(x - y) <= ABS_TOL + REL_TOL * scale;
    }
}
\end{unipdcode}


\begin{example}[Confronto dopo un calcolo approssimato]
Si consideri il quadrato di una radice calcolata in virgola mobile:

\begin{unipdcode}[language=Java]
double r = Math.sqrt(2);
boolean uguali = Numeric.approxEqual(r * r, 2.0);
\end{unipdcode}
Nel calcolo di \(r^2\) può introdursi un errore di
arrotondamento; il controllo con tolleranza valuta
se il risultato sia sufficientemente vicino a \(2\).
I valori \texttt{NaN} e le infinità richiedono una
scelta distinta e non sono considerati approssimativamente
uguali dalla funzione proposta.
\end{example}

\clearpage

\subsection{Metodi di confronto della classe \texttt{Double}}
\label{sec:fi-java-classi-double}

La classe involucro \texttt{Double}, distinta dal tipo
primitivo \texttt{double}, fornisce metodi di confronto.
\texttt{Double.compare(a,b)} restituisce un intero negativo,
nullo o positivo secondo l'ordinamento definito dalla
classe; \texttt{a.equals(b)} confronta invece due oggetti
\texttt{Double} secondo la loro regola di uguaglianza.
Nessuno dei due applica una tolleranza numerica:

\begin{unipdcode}[language=Java]
Double a = 0.1 + 0.2;
Double b = 0.3;
boolean c1 = Double.compare(a, b) == 0; // false
boolean c2 = a.equals(b);               // false
\end{unipdcode}

Le due operazioni distinguono anche \texttt{+0.0}
da \texttt{-0.0} e considerano uguali due valori
\texttt{NaN}, mentre l'operatore \texttt{==}
sui corrispondenti valori primitivi considera uguali
i due zeri e non considera \texttt{NaN} uguale a sé
stesso (\cref{tab:fi-valori-speciali}). Per un confronto
numerico occorre quindi scegliere consapevolmente fra
uguaglianza dei valori rappresentati, ordinamento
definito dalla classe e vicinanza entro una tolleranza
stabilita dal problema.
